\begin{center}
\LARGE
CASE for concurrent systems based on the computer algebra
                         of n-categories

\end{center}

\begin{center}
\Large
Richard Buckland

\end{center}

\noindent
Date: July 18th (Thursday)\\
Time:  14:45-15:05\\

\begin{center}
\large
Abstract
\end{center}

Richard Buckland will describe a concrete application of computational category theory in which
a system designed to perform computations in n-categories is used as the engine for a computer
assisted software engineering (CASE) tool for concurrent systems. The skeletal structure of higher
dimensional categories has been found to provide a useful model of concurrent computation.
However calculating with such higher dimensional algebraic objects is, as is well known to those
who have done such calculations, surprisingly difficult. This difficulty arises both from the
combinatorial complexity of all but the most trivial or low dimensional examples, and from the
problem of verifying the reasonableness of procedures or results since we have little native
intuition when dealing with high dimensions. To assist in this difficult task of high dimensional
calculation a computational system has been developed to perform calculations with Schemes
(schemes are the skeletal structures underlying paths in omega categories). This talk will briefly
describe the algorithms for calculating with Schemes and will then show how, in a nice piece of
circularity, this computer-aided mathematics has as an application the engine of a CASE tool to
specify computational systems. 

